\documentclass[12pt]{article}
\usepackage{mathtools}
\usepackage{amssymb}
\usepackage{fontspec}
\setmainfont{EB Garamond}

% Set page size and margins
\usepackage[a4paper,top=2cm,bottom=1.5cm,left=1cm,right=1cm]{geometry}
\usepackage{fancyhdr}
\pagestyle{fancy}
\setlength{\parindent}{0pt}

% Useful packages
\usepackage{minted}
\usepackage{graphicx}
\graphicspath{{./media/}}
\usepackage{subfig}
\usepackage{float}
\usepackage[colorlinks=true, allcolors=blue]{hyperref}
\usepackage{siunitx}


\title{\vspace{-2cm} Θεμελιώσεις Κρυπτογραφίας \\
        Homework 3}
\author{Μαρινόπουλος Χριστόφορος 10522}
\date{Εαρινό εξάμηνο 2025}

\lhead{Homework 3}
\chead{Θεμελιώσεις Κρυπτογραφίας}
\rhead{Μαρινόπουλος Χριστόφορος}
\begin{document}
\maketitle

\section*{Εκφώνηση}
Δίνεται ένα LFSR μήκους 4 με πολυώνυμο ανάδρασης: \( f(x)=x^4 + x^2 +x + 1 \). Έστω ότι το seed είναι: \( (x[3], x[2], x[1], x[0]) \). Την έξοδο την συμβολίζουμε με: \( y[0], y[1], y[2], ...\).
\subsection*{Ερώτημα 1}
Nα δώσετε τύπους για τα \( y[i] \) για \(i=0, 1, …, 8 \), συναρτήσει των \(x[0], x[1], x[2], x[3] \)
Δηλ. Αν κάποιος έχει τo seed \( (x[3], x[2], x[1], x[0]) \) να μπορεί να υπολογίσει τα
\(y[0], y[1], ..., y[8] \).
\subsection*{Απάντηση}
\begin{align*}
y[0] &= x[3] \\
y[1] &= x[2] \\
y[2] &= x[1] \\
y[3] &= x[0] \\
y[4] &= x[3] \oplus x[1] \oplus x[0] \\
y[5] &= x[2] \oplus x[0] \oplus y[4] \\
             &= x[2] \oplus x[0] \oplus x[3] \oplus x[1] \oplus x[0] \\
y[6] &= x[1] \oplus y[4] \oplus y[5] \\
             &= x[1] \oplus x[3] \oplus x[1] \oplus x[0] \oplus x[2] \oplus x[0] \oplus x[3] \oplus x[1] \oplus x[0] \\
y[7] &= x[0] \oplus y[5] \oplus y[6] \\
             &= x[0] \oplus x[2] \oplus x[0] \oplus x[3] \oplus x[1] \oplus x[0] \oplus x[1] \oplus x[3] \oplus x[1] \oplus x[0] \oplus x[2] \oplus x[0] \oplus x[3] \oplus x[1] \oplus x[0] \\
y[8] &= y[4] \oplus y[6] \oplus y[7] \\
             &= x[3] \oplus x[1] \oplus x[0] \oplus x[1] \oplus x[3] \oplus x[1] \oplus x[0] \oplus x[2] \oplus x[0] \\
             &\quad \oplus x[3] \oplus x[1] \oplus x[0] \oplus x[0] \oplus x[2] \oplus x[0] \oplus x[3] \oplus x[1] \oplus x[0] \\
             &\quad \oplus x[1] \oplus x[3] \oplus x[1] \oplus x[0] \oplus x[2] \oplus x[0] \oplus x[3] \oplus x[1] \oplus x[0]
\end{align*}

Εφόσον \( A \oplus A = 0 \) οι παραπάνω σχέσεις μπορούν να έρθουν στην ακόλουθη μορφή:

\begin{equation*}
\begin{array}{c}
     
\mathbf{y[0] = x[3]} \\
\mathbf{y[1] = x[2]} \\
\mathbf{y[2] = x[1]} \\
\mathbf{y[3] = x[0]} \\
\mathbf{y[4] = x[3] \oplus x[1] \oplus x[0]} \\
\mathbf{y[5] = x[2]  \oplus x[3] \oplus x[1]} \\
\mathbf{y[6] = x[0] \oplus x[2]  \oplus x[1]} \\
\mathbf{y[7] = x[3] }   \\
\mathbf{y[8] = x[2]}    
\end{array}            
\end{equation*}
\hfill \blacksquare
\subsection*{Ερώτημα 2}
Αν \( y[5]=1, y[6]=0, y[7]=0, y[8]=1 \) βρείτε το αρχικό seed (χωρίς brute force).
\subsection*{Απάντηση}
Αξιοποιώντας το προηγούμενο ερώτημα έχουμε:
\begin{equation*}
\begin{array}{c}
y[5] = x[2]  \oplus x[3] \oplus x[1] = 1 \\
y[6] = x[0] \oplus x[2]  \oplus x[1] = 0 \\
y[7] = x[3] = 0 \\
y[8] = x[2] = 1
\end{array}
\end{equation*}
Συνεπώς:
\begin{equation*}
\begin{array}{c}
0 \oplus x[1] = 1 \Rightarrow x[1] = 1 \\
x[0] \oplus 1  \oplus x[1] = 0 \Rightarrow 0\oplus x[0] = 0 \Rightarrow x[0] = 0
\end{array}
\end{equation*}
Επομένως το ζητούμενο seed είναι: \( \mathbf{x[0, 1, 2, 3] = [0, 1, 1, 0]} \)
\hfill \blacksquare
\section*{Bonus}

Να αναπτυχθεί αλγόριθμος (με οποιαδήποτε γλώσσα προγραμματισμού της επιλογής σας) ο 
οποιός θα υπολογίζει ένα πολυώνυμο ανάδρασης, τάξης 64, τέτοιο ώστε ο LFSR που ορίζεται 
να έχει την μέγιστη περίοδο. Επίσης, να δωθεί η πιθάνοτητα επιτυχίας του αλγορίθμου στην πρώτη του προσπάθεια.

\subsection*{Απάντηση}
\[
P(\text{Success}) = \frac{\text{Number of max period 64 bit LFSR}}{\text{Number of every 64 bit LFSR}} = \frac{\frac{\text{φ}(2^{64} - 1)}{64}}{2^{63}} = 1.56%
\]
Όπου \( \text{φ}(x) \) η συνάρτηση του Euler
\hfill \blacksquare

\end{document}